\documentclass[10pt,a4paper]{amsart}
\usepackage{amsmath,amssymb,amsthm,graphicx}
% Retained preamble for the Stacks Project chapter morphisms.tex (Morphisms of Schemes).
% The chapter file pinned-original-source.tex carries only \input{preamble} on its first line, so no
% preamble text lives inside the chapter itself.  The Stacks Project uses one single shared preamble for
% every chapter; the snapshot stages it verbatim in sources/stacks-chow/preamble.tex.  The block below is
% copied mechanically, line for line, from that file's lines 14-18 (the Xy-pic setup) and lines 154-217
% (the native theorem declarations together with the mathematical macros).  Nothing was added, removed or
% redefined.  The wrapper class and the amsmath/amssymb/amsthm/graphicx base packages are supplied by the
% assembler, so the class/\externaldocument part of the shared preamble is not repeated here.
% For commutative diagrams we use Xy-pic
\usepackage[all]{xy}

% We use 2cell for 2-commutative diagrams.
\xyoption{2cell}
% Theorem environments.
%
\theoremstyle{plain}
\newtheorem{theorem}[subsection]{Theorem}
\newtheorem{proposition}[subsection]{Proposition}
\newtheorem{lemma}[subsection]{Lemma}

\theoremstyle{definition}
\newtheorem{definition}[subsection]{Definition}
\newtheorem{example}[subsection]{Example}
\newtheorem{exercise}[subsection]{Exercise}
\newtheorem{situation}[subsection]{Situation}

\theoremstyle{remark}
\newtheorem{remark}[subsection]{Remark}
\newtheorem{remarks}[subsection]{Remarks}

\numberwithin{equation}{subsection}

% Macros
%
\def\lim{\mathop{\mathrm{lim}}\nolimits}
\def\colim{\mathop{\mathrm{colim}}\nolimits}
\def\Spec{\mathop{\mathrm{Spec}}}
\def\Hom{\mathop{\mathrm{Hom}}\nolimits}
\def\Ext{\mathop{\mathrm{Ext}}\nolimits}
\def\SheafHom{\mathop{\mathcal{H}\!\mathit{om}}\nolimits}
\def\SheafExt{\mathop{\mathcal{E}\!\mathit{xt}}\nolimits}
\def\Sch{\mathit{Sch}}
\def\Mor{\mathop{\mathrm{Mor}}\nolimits}
\def\Ob{\mathop{\mathrm{Ob}}\nolimits}
\def\Sh{\mathop{\mathit{Sh}}\nolimits}
\def\NL{\mathop{N\!L}\nolimits}
\def\CH{\mathop{\mathrm{CH}}\nolimits}
\def\proetale{{pro\text{-}\acute{e}tale}}
\def\etale{{\acute{e}tale}}
\def\QCoh{\mathit{QCoh}}
\def\Ker{\mathop{\mathrm{Ker}}}
\def\Im{\mathop{\mathrm{Im}}}
\def\Coker{\mathop{\mathrm{Coker}}}
\def\Coim{\mathop{\mathrm{Coim}}}

% Boxtimes
%
\DeclareMathSymbol{\boxtimes}{\mathbin}{AMSa}{"02}

%
% Macros for moduli stacks/spaces
%
\def\QCohstack{\mathcal{QC}\!\mathit{oh}}
\def\Cohstack{\mathcal{C}\!\mathit{oh}}
\def\Spacesstack{\mathcal{S}\!\mathit{paces}}
\def\Quotfunctor{\mathrm{Quot}}
\def\Hilbfunctor{\mathrm{Hilb}}
\def\Curvesstack{\mathcal{C}\!\mathit{urves}}
\def\Polarizedstack{\mathcal{P}\!\mathit{olarized}}
\def\Complexesstack{\mathcal{C}\!\mathit{omplexes}}
% \Pic is the operator that assigns to X its picard group, usage \Pic(X)
% \Picardstack_{X/B} denotes the Picard stack of X over B
% \Picardfunctor_{X/B} denotes the Picard functor of X over B
\def\Pic{\mathop{\mathrm{Pic}}\nolimits}
\def\Picardstack{\mathcal{P}\!\mathit{ic}}
\def\Picardfunctor{\mathrm{Pic}}
\def\Deformationcategory{\mathcal{D}\!\mathit{ef}}

% --- appended 2026-10-12 (Chief): the project own definitions of the reference/slogan/history
% display environments, copied mechanically from the shared Stacks preamble lines 8-12
% (\usepackage{verbatim} plus the three \newenvironment lines, which the chapter preamble keeps
% verbatim from the project preamble). Nothing added, removed or redefined: these environments
% carry no mathematical content and the project itself typesets them as comments. Without these
% three lines a counted statement whose own text uses them fails with Environment slogan undefined.
\usepackage{verbatim}
\newenvironment{reference}{\comment}{\endcomment}
\newenvironment{slogan}{\comment}{\endcomment}
\newenvironment{history}{\comment}{\endcomment}
\usepackage{cleveref}
\crefname{originallemma}{original lemma}{original lemmas}
\Crefname{originallemma}{Original lemma}{Original lemmas}
\makeatletter
\newlabel{algebra-lemma-catenary}{{02I2}{source}{}{sourcepointer-1}{}}
\newlabel{algebra-lemma-catenary@cref}{{[originallemma][0][]02I2}{[1][0][]source}{}{}{}}
\makeatother
\crefname{originallemma}{original lemma}{original lemmas}
\Crefname{originallemma}{Original lemma}{Original lemmas}
\makeatletter
\newlabel{lemma-composition-finite-type}{{01T3}{source}{}{sourcepointer-2}{}}
\newlabel{lemma-composition-finite-type@cref}{{[originallemma][0][]01T3}{[1][0][]source}{}{}{}}
\makeatother
\crefname{originallemma}{original lemma}{original lemmas}
\Crefname{originallemma}{Original lemma}{Original lemmas}
\makeatletter
\newlabel{lemma-immersion-locally-finite-type}{{01T5}{source}{}{sourcepointer-3}{}}
\newlabel{lemma-immersion-locally-finite-type@cref}{{[originallemma][0][]01T5}{[1][0][]source}{}{}{}}
\makeatother
\crefname{originallemma}{original lemma}{original lemmas}
\Crefname{originallemma}{Original lemma}{Original lemmas}
\makeatletter
\newlabel{lemma-locally-finite-type-characterize}{{01T2}{source}{}{sourcepointer-4}{}}
\newlabel{lemma-locally-finite-type-characterize@cref}{{[originallemma][0][]01T2}{[1][0][]source}{}{}{}}
\makeatother
\crefname{originallemma}{original lemma}{original lemmas}
\Crefname{originallemma}{Original lemma}{Original lemmas}
\makeatletter
\newlabel{properties-lemma-catenary-local}{{02IX}{source}{}{sourcepointer-5}{}}
\newlabel{properties-lemma-catenary-local@cref}{{[originallemma][0][]02IX}{[1][0][]source}{}{}{}}
\makeatother
\title{Universal catenarity of a locally Noetherian scheme is local, and passes to schemes locally of finite type over it}
\author{The Stacks Project Authors}
\date{stacks-project:morphisms.tex:local-snapshot-20261012; GNU FDL 1.2 or later}
\begin{document}
\maketitle
\noindent\textit{Editorial note.} All delivered mathematical text is the author's own, copied from the pinned Stacks Project chapter file pinned-original-source.tex; no statement or proof step was rewritten, reordered or removed.\par\medskip
\noindent\textit{Editorial note.} Following the multi-claim convention for this source, source\_url is the source-level URL https://stacks.math.columbia.edu/ and the source label lemma-universally-catenary-local is carried in claim\_key and here rather than in a per-claim tag URL.\par\medskip
\noindent\textit{Editorial note.} The preamble is the chapter's shared Stacks Project preamble, retained as sources/stacks-morphisms/retained-preamble.tex: a line-for-line excerpt of the shared Stacks preamble staged with this snapshot. No macro was invented or redefined, and this card adds no preamble text of its own.\par\medskip
\noindent\textit{Editorial note.} Statement span 3118-3132 and proof span 3134-3161 come from the halfA pool card operations/pool-cards/stacks-morphisms-halfA.json and were re-read in the pinned chapter file before this card was written.\par\medskip
\noindent\textit{Editorial note.} The proof refers to the statement's conditions by numbers (1) to (4); those references are the author's and are preserved.\par\medskip
\noindent\textit{Editorial note.} 5 cross-references occur inside the counted spans; they are listed in external\_references as navigation pointers only.\par\medskip
\noindent\textit{Allowed standard prerequisites (editorial).}\begin{itemize}
\item An open immersion is locally of finite type and a composite of morphisms locally of finite type is locally of finite type, both cited from the same chapter. Cited at proof lines 3351-3352. Shows that open subschemes and schemes locally of finite type over a universally catenary base are again universally catenary.
\item The algebra lemma that over a universally catenary ring every finite type algebra is catenary. Cited at proof lines 3360-3361. Converts the hypothesis on one affine chart into universal catenarity of its coordinate ring.
\item The same-chapter characterization of locally finite type morphisms on affine charts, and the properties chapter lemma that catenarity is local on a scheme. Cited at proof lines 3372-3373 and 3376-3377. Make the coordinate rings of the chosen affine cover of X finite type algebras over the base rings, and then assemble catenarity from the pieces.
\end{itemize}\medskip
\ifdefined\hypertarget\hypertarget{sourcepointer-1}{}\fi
\noindent\textit{Source reference.} Original lemma 02I2: focused-seeds/AGP/sources/stacks-snapshot-20261011/tags.txt, source lines 3285--3285. Same-chapter reference used in the retained author text at source line(s) 3147. It is supplied only as a navigation pointer so that the retained text compiles; the counted item delivers the author's own statement and the author's complete proof, and no part of the referenced item is copied in or substituted for the counted argument.\par\medskip
\ifdefined\hypertarget\hypertarget{sourcepointer-2}{}\fi
\noindent\textit{Source reference.} Original lemma 01T3: focused-seeds/AGP/sources/stacks-snapshot-20261011/tags.txt, source lines 2360--2360. Same-chapter reference used in the retained author text at source line(s) 2889, 2937, 3138, 3329, 9763, 9913. It is supplied only as a navigation pointer so that the retained text compiles; the counted item delivers the author's own statement and the author's complete proof, and no part of the referenced item is copied in or substituted for the counted argument.\par\medskip
\ifdefined\hypertarget\hypertarget{sourcepointer-3}{}\fi
\noindent\textit{Source reference.} Original lemma 01T5: focused-seeds/AGP/sources/stacks-snapshot-20261011/tags.txt, source lines 2362--2362. Same-chapter reference used in the retained author text at source line(s) 2980, 3135, 3326, 9944. It is supplied only as a navigation pointer so that the retained text compiles; the counted item delivers the author's own statement and the author's complete proof, and no part of the referenced item is copied in or substituted for the counted argument.\par\medskip
\ifdefined\hypertarget\hypertarget{sourcepointer-4}{}\fi
\noindent\textit{Source reference.} Original lemma 01T2: focused-seeds/AGP/sources/stacks-snapshot-20261011/tags.txt, source lines 2359--2359. Same-chapter reference used in the retained author text at source line(s) 2717, 2739, 2801, 2863, 3019, 3035, 3157, 3348, 3613, 3707, 3868, 5197, 5282, 6702, 9325, 9390, 9440, 11421, 13422, 13686. It is supplied only as a navigation pointer so that the retained text compiles; the counted item delivers the author's own statement and the author's complete proof, and no part of the referenced item is copied in or substituted for the counted argument.\par\medskip
\ifdefined\hypertarget\hypertarget{sourcepointer-5}{}\fi
\noindent\textit{Source reference.} Original lemma 02IX: focused-seeds/AGP/sources/stacks-snapshot-20261011/tags.txt, source lines 3316--3316. Same-chapter reference used in the retained author text at source line(s) 3160. It is supplied only as a navigation pointer so that the retained text compiles; the counted item delivers the author's own statement and the author's complete proof, and no part of the referenced item is copied in or substituted for the counted argument.\par\medskip
\begin{lemma}
\label{lemma-universally-catenary-local}
Let $S$ be a locally Noetherian scheme. The following are equivalent
\begin{enumerate}
\item $S$ is universally catenary,
\item there exists an open covering of $S$ all of whose members are
universally catenary schemes,
\item for every affine open $\Spec(R) = U \subset S$ the ring
$R$ is universally catenary, and
\item there exists an affine open covering $S = \bigcup U_i$ such
that each $U_i$ is the spectrum of a universally catenary ring.
\end{enumerate}
Moreover, in this case any scheme locally of finite type over $S$
is universally catenary as well.
\end{lemma}

\begin{proof}
By Lemma \ref{lemma-immersion-locally-finite-type} an open immersion
is locally of finite type. A composition of morphisms locally of
finite type is locally of finite type
(Lemma \ref{lemma-composition-finite-type}). Thus it is clear that if $S$ is
universally catenary then any open and any scheme locally of finite
type over $S$ is universally catenary as well. This proves the final
statement of the lemma and that (1) implies (2).

\medskip\noindent
If $\Spec(R)$ is a universally catenary scheme, then every
scheme $\Spec(A)$ with $A$ a finite type $R$-algebra is
catenary. Hence all these rings $A$ are catenary by
Algebra, Lemma \ref{algebra-lemma-catenary}.
Thus $R$ is universally catenary. Combined with the remarks above we
conclude that (1) implies (3), and (2) implies (4). Of course
(3) implies (4) trivially.

\medskip\noindent
To finish the proof we show that (4) implies (1).
Assume (4) and let $X \to S$ be a morphism locally of finite type.
We can find an affine open covering $X = \bigcup V_j$ such that
each $V_j \to S$ maps into one of the $U_i$. By
Lemma \ref{lemma-locally-finite-type-characterize}
the induced ring map $\mathcal{O}(U_i) \to \mathcal{O}(V_j)$ is
of finite type. Hence $\mathcal{O}(V_j)$ is catenary. Hence
$X$ is catenary by Properties, Lemma \ref{properties-lemma-catenary-local}.
\end{proof}

\end{document}
